\providecommand{\kch}{\kappa_{\mathrm{ch}}}
\providecommand{\kcat}{\kappa_{\mathrm{cat}}}
\providecommand{\kBKM}{\kappa_{\mathrm{BKM}}}
\providecommand{\kfib}{\kappa_{\mathrm{fiber}}}
\providecommand{\Meff}{\mathcal{M}^{+}_{\mathrm{eff}}}
\providecommand{\Mred}{\mathcal{M}^{\mathrm{red}}_{\mathrm{DT}}}
\providecommand{\virred}{[\Mred]^{\mathrm{vir,red}}}
\providecommand{\Etr}{E_{\mathrm{tr}}}
\providecommand{\Trans}{\mathrm{Trans}}
\providecommand{\PhiTenOP}{\Phi_{10}^{\mathrm{OP}}}
\providecommand{\PhiTenUn}{\Phi_{10}^{\mathrm{un}}}
\providecommand{\DeltaFive}{\Delta_{5}}
\providecommand{\DFive}{D_{5}}
\providecommand{\gDeltaFive}{\mathfrak{g}_{\DeltaFive}}
\providecommand{\gBPS}{\mathfrak{g}_{\mathrm{BPS}}}
\providecommand{\Yplus}{Y^{+}}
\providecommand{\CoHAred}{\mathrm{CoHA}^{\mathrm{red}}}
\providecommand{\K}{K3}
\providecommand{\Db}{D^b}

\section*{The \(N=1\) effective positive half on \(\K\times E\)}
\label{sec:n1-effective-positive-half}

\subsection*{Primitive reduced sector}

Let \(X=\K\times E\). The \(N=1\) sector is the identity-twined
primitive sector. A charge is written
\[
  \gamma=(r,\beta_h,n;\ell),
  \qquad
  \beta_h\in H_2(\K,\mathbb Z),\quad
  \beta_h^2=2h-2,\quad
  \ell\in H_2(E,\mathbb Z)\cong\mathbb Z,
\]
with \(\beta_h\) primitive. The pullback of the holomorphic symplectic
form on \(\K\) gives the semiregularity cosection
\[
  \mathrm{Ob}\twoheadrightarrow\mathcal O .
\]
The reduced obstruction theory is obtained by deleting this trivial
quotient:
\[
  \mathrm{Ob}^{\mathrm{red}}
  =
  \ker(\mathrm{Ob}\twoheadrightarrow\mathcal O).
\]
For primitive \(\K\)-classes, Maulik--Pandharipande--Thomas and
Oberdieck--Pandharipande construct the corresponding reduced virtual
class. The translation group
\[
  \Etr=\Trans(E)\cong E
\]
acts on the elliptic factor. The quotient stack carries the
zero-dimensional reduced class used in the primitive reduced
Donaldson--Thomas integral:
\[
  \Mred(X;\gamma)
  =
  \bigl\{I^\bullet\in\Db(X)\mid
  \operatorname{ch}(I^\bullet)=\gamma,\ \beta_h\ \mathrm{primitive}\bigr\},
  \qquad
  \virred\in A_0\bigl(\Mred(X;\gamma)/\Etr\bigr).
\]
Kawai--Yoshioka give
\[
  \sum_{n\geq 0}\chi(\operatorname{Hilb}^n\K)q^n
  =
  \prod_{m\geq 1}(1-q^m)^{-24}
  =
  q\,\eta(q)^{-24}.
\]
Equivalently, after the standard reduced-DT shift, the seed is
\(\eta(q)^{-24}\).

\subsection*{Effective Hall positive half}

Fix a stability chamber, algebraic orientation square roots for the
reduced \(d\)-critical stack, and compact supports for the Hall
correspondences. The effective primitive stack is
\[
  \Meff(X)\big|_{N=1}
  =
  \bigsqcup_{\gamma\in\mathrm{Eff}^{\mathrm{prim}}_{\K}}
  \Mred(X;\gamma)/\Etr .
\]
The oriented effective positive half is
\[
  \Yplus(X)\big|_{N=1}
  =
  \bigoplus_{\gamma\in \mathrm{Eff}^{\mathrm{prim}}_{\K}}
  H^{\mathrm{BM}}_\bullet\!\left(
    \Mred(X;\gamma)/\Etr,\,
    \phi_{W_\gamma}\otimes\mathscr L_\gamma^{1/2}
  \right),
\]
where \(\phi_{W_\gamma}\) is the Kontsevich--Soibelman
vanishing-cycle sheaf on local critical charts and
\(\mathscr L_\gamma^{1/2}\) is the chosen orientation square root. Hall
convolution along the effective cone makes this an \(E_1\) algebra. In
complex dimension \(3\), the \(E_2\) structure belongs to centre,
double, or module layers built from the \(E_1\) algebra.

\subsection*{Oberdieck--Pixton scalar}

The reduced DT scalar in the Oberdieck--Pixton monic normalisation is
\begin{equation}
\label{eq:n1-op-scalar}
  Z^{\mathrm{red}}_{\mathrm{DT}}(X;q,t,p)
  =
  -\bigl(\PhiTenOP(q,t,p)\bigr)^{-1}
  =
  -\DFive(q,t,p)^{-2}
  =
  -4096\,\DeltaFive(q,t,p)^{-2}.
\end{equation}
The variables are
\[
  q=e^{2\pi i\tau},\qquad
  t=e^{2\pi iz},\qquad
  p=e^{2\pi i\sigma}.
\]
The normalisations are
\[
  \DFive=64^{-1}\DeltaFive,\qquad
  \PhiTenUn=\DeltaFive^2,\qquad
  \PhiTenOP=\DFive^2=4096^{-1}\DeltaFive^2.
\]
Thus \(4096=64^2\) is the theta-normalised to monic conversion squared.
The Behrend sign is the leading minus sign in
\eqref{eq:n1-op-scalar}.

The monic primitive denominator is
\[
  \DFive(q,t,p)
  =
  q^{1/2}t^{1/2}p^{1/2}
  \prod_{\substack{(n,\ell,m)>0\\4nm-\ell^2\ge -1}}
  \bigl(1-q^n t^\ell p^m\bigr)^{c_1(4nm-\ell^2)}.
\]
Consequently
\[
  \PhiTenOP(q,t,p)
  =
  qtp
  \prod_{\substack{(n,\ell,m)>0\\4nm-\ell^2\ge -1}}
  \bigl(1-q^n t^\ell p^m\bigr)^{2c_1(4nm-\ell^2)}.
\]
Here
\[
  \phi_{0,1}(\tau,z)
  =
  \sum_{n,\ell}c_1(4n-\ell^2)q^n t^\ell,
  \qquad
  c_1(0)=10.
\]
The scalar \eqref{eq:n1-op-scalar} is the signed reduced DT character.
It is separate from the effective Hall positive half and separate from
the Hall--Drinfeld double.

\subsection*{Primitive denominator and Hall recognition}

The primitive \(N=1\) Borcherds denominator is \(\DeltaFive\), or
equivalently its monic form \(\DFive=64^{-1}\DeltaFive\). Scalar
rescaling changes neither divisor nor weight:
\[
  \kBKM(\DeltaFive)=\kBKM(\DFive)=\frac{c_1(0)}{2}=5.
\]
The Gritsenko--Nikulin Borcherds--Kac--Moody target has signed root
multiplicities
\begin{equation}
\label{eq:n1-target-multiplicity}
  \operatorname{sdim}(\gDeltaFive)_{\alpha(n,\ell,m)}
  =
  c_1(4nm-\ell^2).
\end{equation}
The OP scalar has squared exponents \(2c_1(4nm-\ell^2)\). The primitive
BKM target uses the square root denominator \(\DeltaFive\), hence the
unsquared exponents in \eqref{eq:n1-target-multiplicity}.

Let
\[
  \gBPS(X):=\operatorname{Prim}\,\CoHAred(X)
\]
when the reduced critical CoHA, its orientation, and compact-support
functoriality have been constructed. A source-side identification
\[
  \operatorname{sdim}\gBPS(X)_{(n,\ell,m)}
  =
  \operatorname{sdim}(\gDeltaFive)_{\alpha(n,\ell,m)}
\]
requires the Hall--Drinfeld recognition data: a Serre-dual negative
half, a non-degenerate Hopf pairing after radical quotient, a
coproduct, completion, centre compatibility, parity control, and a PBW
comparison with the BKM target. The scalar identity
\eqref{eq:n1-op-scalar} is a required numerical test for these data.

\subsection*{Automorphic closure}

Borcherds' singular theta lift gives the automorphic product
\[
  \operatorname{Bor}^{\mathrm{mon}}(\phi_{0,1})=\DFive,\qquad
  \operatorname{Bor}^{\mathrm{mon}}(2\phi_{0,1})=\PhiTenOP=\DFive^2.
\]
In the theta-normalised convention this reads
\[
  \operatorname{Bor}^{\theta}(\phi_{0,1})=\DeltaFive,\qquad
  \PhiTenUn=\DeltaFive^2.
\]
The weight formula gives
\[
  \operatorname{wt}(\DeltaFive)
  =
  \operatorname{wt}(\DFive)
  =
  \frac{c_1(0)}{2}
  =
  5.
\]
The square \(\PhiTenUn=\DeltaFive^2\) has weight \(10\), and
\(\PhiTenOP=\DFive^2\) has the same weight. The BKM invariant in this
sector belongs to the primitive denominator \(\DeltaFive\); the squared
scalar is the OP/Igusa character layer.

\begin{theorem}[\(N=1\) effective positive half and scalar test]
\label{thm:n1-effective-positive-half}
For \(X=\K\times E\) in the identity-twined primitive sector:
\begin{enumerate}
\item
\[
  \Meff(X)\big|_{N=1}
  =
  \bigsqcup_{\gamma\in\mathrm{Eff}^{\mathrm{prim}}_{\K}}
  \Mred(X;\gamma)/\Etr,
  \qquad
  \Etr=\Trans(E)\cong E .
\]
\item
\(\Yplus(X)|_{N=1}\) is the oriented Borel--Moore
vanishing-cycle homology of this effective stack, with \(E_1\) Hall
convolution along the effective cone.
\item
The reduced DT scalar is
\[
  Z^{\mathrm{red}}_{\mathrm{DT}}(X)
  =
  -(\PhiTenOP)^{-1}
  =
  -\DFive^{-2}
  =
  -4096\,\DeltaFive^{-2}.
\]
\item
The primitive denominator satisfies
\[
  \DFive=64^{-1}\DeltaFive,\qquad
  \PhiTenOP=\DFive^2=4096^{-1}\DeltaFive^2,\qquad
  \kBKM(\DeltaFive)=\frac{c_1(0)}{2}=5.
\]
\item
The target BKM signed root multiplicities are
\[
  \operatorname{sdim}(\gDeltaFive)_{\alpha(n,\ell,m)}
  =
  c_1(4nm-\ell^2).
\]
Their realisation by \(\operatorname{Prim}\CoHAred(X)\) requires the
Hall--Drinfeld recognition data listed above.
\end{enumerate}
\end{theorem}

\begin{proof}
The K3 symplectic form gives the cosection
\(\mathrm{Ob}\twoheadrightarrow\mathcal O\); its kernel is the reduced
obstruction theory. The elliptic translation group acts on the \(E\)
factor, and the quotient stack \(\Mred(X;\gamma)/\Etr\) is the
zero-dimensional reduced sector used in the primitive integral.
Orientation square roots and compact-support Hall correspondences make
the direct sum of Borel--Moore vanishing-cycle groups into the
effective \(E_1\) Hall positive half.

Oberdieck--Pandharipande's reduced DT theorem, read in the
Oberdieck--Pixton monic convention, gives
\(-(\PhiTenOP)^{-1}\). Gritsenko--Nikulin's theta-normalised form has
leading coefficient \(64\), so \(\DFive=64^{-1}\DeltaFive\) and
\(\PhiTenOP=\DFive^2=4096^{-1}\DeltaFive^2\). The product for
\(\DFive\) has exponents \(c_1(4nm-\ell^2)\); squaring gives the OP
exponents \(2c_1(4nm-\ell^2)\). Borcherds' weight formula gives
\(\kBKM(\DeltaFive)=c_1(0)/2=10/2=5\). The signed root multiplicities
of the primitive BKM target are therefore the unsquared exponents.
\end{proof}

\subsection*{Subscripted invariants on \(K3\times E\)}

The compact total-space categorical value is K\"unneth-multiplicative:
\[
  \kcat(X)
  =
  \chi(\mathcal O_X)
  =
  \chi(\mathcal O_{\K})\chi(\mathcal O_E)
  =
  2\cdot 0
  =
  0.
\]
The compact chiral Hodge supertrace is
\[
  \kch(X)
  =
  \sum_q(-1)^q h^{0,q}(X)
  =
  1-1+1-1
  =
  0.
\]
The Heisenberg--Mukai specialisation is a different construction:
\[
  \kch^{\mathrm{Heis}}(X)=3.
\]
The primitive Borcherds denominator gives
\[
  \kBKM(\DeltaFive)=5,
\]
and the K3 fibre lattice gives
\[
  \kfib(X)=\operatorname{rk}H^\bullet(\K,\mathbb Z)=24.
\]
Thus the \(K3\times E\) constants in this sector are
\[
  \bigl(
    \kcat(X),\,
    \kch(X),\,
    \kch^{\mathrm{Heis}}(X),\,
    \kBKM(\DeltaFive),\,
    \kfib(X)
  \bigr)
  =
  (0,0,3,5,24).
\]
These entries come from distinct constructions. The BKM entry is fixed
by \(c_1(0)/2\); the compact entries are fixed by K\"unneth and Hodge
theory; the fibre entry is the Mukai rank.
